\section{Step 71 Separate byte capacity from bit position}\label{h263-writer-s01}
Load the three context doublewords and keep their units distinct. Check context
alignment and its non-wrapping extent before reading fields. Validate capacity,
buffer extent and used-bit range before deriving an output byte address. Use
full-width address arithmetic for the buffer plus byte index.

\paragraph{Observe} An eight-byte arena has sixty-four bit positions. A count
of eight means one byte is committed, not the entire arena.
\Needspace{8\baselineskip}
\paragraph{Check} What unit does the mutable context counter use?
\begin{itemize}
  \item[A.] Bits already committed to the stream.
  \item[B.] Capacity measured in native words.
  \item[C.] Video pictures per second.
\end{itemize}

\section{Step 72 Validate the field value and width}\label{h263-writer-s02}
Reject a width greater than 32 before using it as a variable shift count. Then
check that shifting the unsigned value right by the declared width leaves zero.
For width zero, this checks the complete value directly. Keep both arguments
full width; a large value or width can have a small plausible low word.

\paragraph{Observe} Value 256 does not fit a declared width of eight. Value zero
with width zero is a legal empty field; value one with width zero is invalid.
\Needspace{8\baselineskip}
\paragraph{Check} Why reject high bits instead of silently masking them?
\begin{itemize}
  \item[A.] It changes the memory byte order.
  \item[B.] It catches a malformed field before corrupting its intended syntax.
  \item[C.] It permits widths greater than 32.
\end{itemize}

\section{Step 73 Admit the complete field before touching output}\label{h263-writer-s03}
After checking the existing used count, subtract it from eight times capacity.
Compare the complete requested width with that remaining count. Return status
two before buffer reads or writes if the field does not fit. Preserve all context
fields on rejection. A valid empty append needs no output access or counter store.

\paragraph{Observe} With one byte of capacity and seven used bits, a one-bit
append fits exactly. A two-bit append retains the old byte and old count.
\Needspace{8\baselineskip}
\paragraph{Check} What happens when a valid field exceeds remaining space?
\begin{itemize}
  \item[A.] Write the fitting prefix and return success.
  \item[B.] Advance the count without writing bytes.
  \item[C.] Return status two without accessing the buffer.
\end{itemize}

\section{Step 74 Emit the highest selected bit first}\label{h263-writer-s04}
For the initial scalar loop, begin with source shift $n-1$. Extract one unsigned
bit, derive destination byte $p/8$ and bit position $7-(p\bmod8)$, and place it
there. Advance the local bit cursor and move to the next lower source bit until
all requested bits are consumed.

The initial empty-width return must precede the loop, so there is no source
shift of minus one. The first emitted bit belongs to the declared field, which
can contain leading zeros even when the numerical value is small.

\paragraph{Observe} Value five with width five emits \texttt{00101}, not just
\texttt{101} and not the eight-byte native representation of five.
\Needspace{8\baselineskip}
\paragraph{Check} Which source bit is emitted first for width $n>0$?
\begin{itemize}
  \item[A.] Bit $n-1$ of the right-aligned value.
  \item[B.] Always bit zero.
  \item[C.] The highest bit of the native address.
\end{itemize}

\section{Step 75 Replace the selected destination bits}\label{h263-writer-s05}
Load the touched byte unsigned. Clear the selected destination bit with an
inverted mask, then OR in the shifted replacement bit. Store a byte. This handles
both replacement values while preserving the old stream prefix and unused suffix.
Do not clear the entire byte when appending inside it.

\paragraph{Observe} Writing a zero over an old one changes that bit. OR without
clearing would leave the old one in place.
\Needspace{8\baselineskip}
\paragraph{Check} Why is an OR-only append incorrect?
\begin{itemize}
  \item[A.] It cannot write a one.
  \item[B.] It cannot clear an old one when the new bit is zero.
  \item[C.] It always reverses native byte order.
\end{itemize}

\section{Step 76 Commit the cursor after successful output}\label{h263-writer-s06}
Keep the advancing cursor local until all field bits have been stored. Then write
it to context offset sixteen and return the appended width. Keep the buffer
pointer and byte capacity unchanged. No counter write is needed for an empty
append.

This ordering is a boundary for the admitted operation. It does not promise
recovery from an asynchronous simulator interruption or an invalid unmapped
pointer that violated caller obligations.

\paragraph{Observe} After a successful nine-bit append at used count three,
the context count is twelve and the returned append count is nine.
\Needspace{8\baselineskip}
\paragraph{Check} Which quantity belongs in the successful return count?
\begin{itemize}
  \item[A.] Byte capacity.
  \item[B.] The numerical field value.
  \item[C.] The declared width just appended.
\end{itemize}

\section{Step 77 Observe repeated appends through the caller}\label{h263-writer-s07}
Start with \path{construction/h263-bit-writer-construction/main.asm}. Its first call
appends value five in three bits, and its second appends decimal nineteen in five
bits to the same context. The complete byte is \texttt{10110011}, hexadecimal
\texttt{B3}, or decimal 179. The count is eight.

At \texttt{writer\_caller\_return}, before its first instruction, the second
append has completed. Return status is zero and appended count five, while the
named status/count fields have not yet received their copies. The ordinary caller
prints 179 and exits zero after resuming. An unfinished entry reports development
status four and exits one.

\paragraph{Observe} The second call continues from bit three; it does not
replace the entire first byte or start a new byte automatically.
\Needspace{8\baselineskip}
\paragraph{Check} Which byte results from fields \texttt{101} and \texttt{10011}?
\begin{itemize}
  \item[A.] \texttt{10110011}.
  \item[B.] \texttt{00000101} only.
  \item[C.] Two unrelated native doublewords.
\end{itemize}

\section{Step 78 Combine the next byte fragment}\label{h263-writer-s08}
Replace the single-bit loop with a byte-fragment loop while retaining the same
contract. Take at most the remaining positions in the destination byte. Extract
that many next high source bits, position them with a mask, and perform one
read-modify-write. Advance the local cursor by the fragment width.

For an unaligned 32-bit field, five fragments suffice. Preserve the old suffix
of the final byte. A whole-word load extending beyond the touched byte range
would violate this exercise's access contract even if it looked faster.

\paragraph{Observe} Both schedules produce identical stream bytes. The wider
fragment reduces repeated byte accesses without changing field order.
\Needspace{8\baselineskip}
\paragraph{Check} How many bytes can an unaligned 32-bit append touch?
\begin{itemize}
  \item[A.] Always exactly four.
  \item[B.] At most five.
  \item[C.] Thirty-two native doublewords.
\end{itemize}

\section{Step 79 Check boundaries before comparing work}\label{h263-writer-s09}
Run public feedback on your current source:
\begin{verbatim}
python3 construction/h263-bit-writer-construction/grade.py \
    my-writer.asm --cli /path/to/mars-cli.jar \
    --report writer-results.json
\end{verbatim}
The checker constructs an independent bit sequence and compares the affected
bytes. It checks both simulator byte orders, every starting bit offset, representative
widths including zero and 32, exact fits, insufficient capacity, wide malformed
controls, guarded context and the maximum-capacity endpoint. Read/write traces
reject access outside bytes touched by the field.

Measure entry-to-return instruction work for the same value, width and initial
cursor in both implementations. Include validation, loop branches, byte accesses
and the committed count update. This local work comparison is not a whole-encoder
wall-time measurement or a grading threshold.

For the named all-one 32-bit field starting at bit three, the tested bit-at-a-time
entry executes 656 instructions and the byte-fragment entry executes 189, including
admission and the final counter store. Both simulator byte orders give the same
counts. This saves 467 instructions for that append; it does not measure an entire
coefficient event, file write or encoder launch. Different valid implementations
need not reproduce these counts.

\paragraph{Observe} A bit-at-a-time implementation remains valid after the
byte-fragment implementation is introduced. Correct output does not require a
particular instruction count or private label.
\Needspace{8\baselineskip}
\paragraph{Check} What must match when comparing the two schedules?
\begin{itemize}
  \item[A.] Only their scratch-register names.
  \item[B.] The host computer's display theme.
  \item[C.] Field value, width, initial bit cursor and complete measured boundary.
\end{itemize}

\section{Step 80 Append alignment bits explicitly}\label{h263-writer-s10}
Prepare a crossing observation without changing your instructions:
\begin{verbatim}
python3 construction/h263-bit-writer-construction/prepare.py \
    --source my-writer.asm --first 5 --first-width 3 \
    --second 419 --second-width 9 --output crossing.asm
\end{verbatim}
The two fields commit twelve bits and leave bytes \texttt{BA 35}; the low four
bits are retained unused storage. To reach the next byte boundary where the
syntax requires zero stuffing, append value zero with width four. This produces
\texttt{BA 30} and count sixteen. In general, the required zero-bit count is
$(8-(p\bmod8))\bmod8$. An already aligned stream needs zero additional bits.

Keep alignment decisions in the syntax writer. Picture start codes require their
specified alignment, but ordinary adjacent fields must not acquire arbitrary
padding. The next coefficient and macroblock stages will use this append entry
with their own lengths, admission and publication rules.

\paragraph{Observe} Capacity and unused storage do not decide syntax. Explicit
padding changes both the byte bits and the committed count.
\Needspace{8\baselineskip}
\paragraph{Check} When should ordinary field appends insert padding automatically?
\begin{itemize}
  \item[A.] Never; the syntax layer requests any required alignment bits.
  \item[B.] After every field regardless of its syntax.
  \item[C.] Whenever an unused buffer byte is nonzero.
\end{itemize}
